\subsection*{Validación en vivo del anclaje y la contextualización}

La contextualización generativa se evaluó mediante una campaña en vivo contra
la API oficial de Mistral. La campaña realizó 16 llamadas, una por cada tipo de
ataque reconocido por la taxonomía
\texttt{multidataset\_attack\_type\_v1\_2026-09-06}, y utilizó el catálogo de
amenazas v2.0. El modelo generativo fue
\texttt{ministral-8b-2512}, con temperatura 0 y salida JSON. Antes de iniciar
la campaña se comprobó la disponibilidad del proveedor. El modelo inicialmente
previsto, \texttt{mistral-small-2603}, respondió con limitación de tasa; por
ello se seleccionó una versión fechada disponible y se registró su identificador
exacto. El intento fallido de comprobación no se incluyó entre los 16 casos de
evaluación.

El objetivo de esta campaña fue aislar el tramo de mitigación. No se enviaron
nuevamente entradas crudas al estandarizador: se reutilizaron eventos del nuevo
\textit{split} de prueba que ya habían sido estandarizados mediante Mistral. El
flujo ejecutado fue: paso directo del evento preestandarizado, detector,
clasificador de 16 tipos activado expresamente para la campaña, recuperación
del catálogo, contextualización real con Mistral, juez, persistencia y auditor.
Por tanto, se trata de un \emph{replay} operacional posterior a la
estandarización, no de una prueba completa del \textit{endpoint} público desde
una entrada cruda.

La selección se congeló antes de la primera llamada al proveedor. De las 1.200
filas del \textit{split} de prueba, 785 cumplían simultáneamente las siguientes
condiciones: mapeo exacto, confianza de mapeo mayor o igual que 0,90,
probabilidad del detector mayor o igual que 0,80 y predicción correcta del tipo
con confianza mayor o igual que 0,80. Se escogió un caso por tipo, priorizando
la confianza del clasificador y, después, la del detector; una función SHA-256
con semilla 42 se empleó únicamente para desempatar. La confianza mínima del
clasificador entre los casos finalmente seleccionados fue 0,9947 y la
probabilidad mínima del detector fue 0,9646.

Antes del \emph{replay} se volvió a aplicar el saneamiento contra fuga del
objetivo. Los 16 eventos superaron el contrato que exige que
\texttt{label\_raw}, \texttt{attack\_family} y
\texttt{attack\_subtype} no estén informados. El tipo esperado se usó para
seleccionar y evaluar el caso, pero no se entregó al flujo; el mitigador recibió
la predicción producida por el clasificador. Este saneamiento no equivale a una
anonimización: las direcciones IP, los puertos y la evidencia técnica necesaria
sí se enviaron al proveedor externo.

El modelo recibió las medidas numeradas y las referencias autorizadas para el
tipo predicho. Para asociar una contextualización con el catálogo debía indicar
un \texttt{base\_id} válido y no introducir identificadores ATT\&CK, CAPEC o
M-* ajenos a la entrada recuperada. La acción operativa almacenada en
\texttt{text} permaneció idéntica a la medida catalogada; la redacción del LLM
se conservó separadamente en \texttt{context}, con
\texttt{context\_trusted=false}.

La Tabla~\ref{tab:mitigador_mistral_global} resume los resultados. Las 16
peticiones devolvieron HTTP~200 y un objeto válido conforme al esquema; no se
activó la ruta de reserva. Mistral contextualizó 111 de las 113 medidas base
presentadas (98,23\,\%). Las dos medidas no contextualizadas permanecieron como
literales del catálogo. Además, generó nueve elementos sin correspondencia
catalogada en ocho casos. Estos nueve elementos representan el 7,50\,\% de los
120 textos generados que terminaron clasificados como contextualización
anclada o como sugerencia del LLM.

\begin{table}[H]
    \ttabbox
    {\caption{Resultado global de la campaña en vivo del mitigador}
    \label{tab:mitigador_mistral_global}}
    {\begin{tabular}{P{8.0cm} r}
        \hline
        \textbf{Indicador} & \textbf{Resultado} \\
        \hline
        Llamadas Mistral con HTTP~200 & 16/16 \\
        \hline
        Respuestas válidas según el esquema & 16/16 \\
        \hline
        Casos con \textit{fallback} & 0/16 \\
        \hline
        Casos con las cinco primeras recomendaciones ancladas & 16/16 \\
        \hline
        Medidas base contextualizadas & 111/113 (98,23\,\%) \\
        \hline
        Casos con cobertura completa de las medidas base & 15/16 \\
        \hline
        Sugerencias no catalogadas & 9 en 8 casos \\
        \hline
        Solicitudes de revisión emitidas por la respuesta cruda & 16/16 \\
        \hline
        Solicitudes neutralizadas por la regla de las cinco primeras & 15/16 \\
        \hline
        Decisiones del juez & 15 aprobaciones; 1 revisión \\
        \hline
        Veredictos del auditor & 15 \texttt{approve}; 1 \texttt{review} \\
        \hline
        Fallos duros del auditor & 0 \\
        \hline
        Latencia Mistral mínima / mediana / P95 / máxima &
        8,72 / 10,85 / 14,41 / 16,34 s \\
        \hline
    \end{tabular}}
\end{table}

La Tabla~\ref{tab:mitigador_mistral_tipo} muestra el desglose. La columna
``contextualizadas/base'' indica cuántas medidas catalogadas recibieron un
contexto válido respecto de cuántas se presentaron al modelo. En todos los
tipos, las cinco primeras recomendaciones se asociaron con cinco bases válidas
y distintas.

\begin{table}[H]
    \centering
    \scriptsize
    \caption{Resultado de la contextualización en vivo por tipo de ataque}
    \label{tab:mitigador_mistral_tipo}
    \begin{tabular}{P{3.2cm} P{4.0cm} c c c P{1.8cm}}
        \hline
        \textbf{Tipo} & \textbf{Origen / perfil} &
        \textbf{Context./base} & \textbf{Primeras 5} &
        \textbf{Sug.} & \textbf{Juez} \\
        \hline
        \texttt{Backdoor} & Edge-IIoTset / \texttt{network\_packet} & 7/7 & sí & 0 & aprueba \\
        \hline
        \texttt{DDoS\_HTTP} & Edge-IIoTset / \texttt{network\_packet} & 7/7 & sí & 1 & aprueba \\
        \hline
        \texttt{DDoS\_ICMP} & Edge-IIoTset / \texttt{network\_packet} & 7/7 & sí & 1 & aprueba \\
        \hline
        \texttt{DDoS\_TCP} & Edge-IIoTset / \texttt{network\_packet} & 7/7 & sí & 1 & aprueba \\
        \hline
        \texttt{DDoS\_UDP} & Edge-IIoTset / \texttt{network\_packet} & 5/7 & sí & 2 & revisión \\
        \hline
        \texttt{Fingerprinting} & Edge-IIoTset / \texttt{network\_packet} & 7/7 & sí & 1 & aprueba \\
        \hline
        \texttt{MITM} & TON-IoT/red / \texttt{network\_flow} & 7/7 & sí & 1 & aprueba \\
        \hline
        \texttt{Password} & TON-IoT/red / \texttt{network\_flow} & 7/7 & sí & 0 & aprueba \\
        \hline
        \texttt{Port\_Scanning} & Edge-IIoTset / \texttt{network\_packet} & 7/7 & sí & 0 & aprueba \\
        \hline
        \texttt{Ransomware} & TON-IoT/red / \texttt{network\_flow} & 7/7 & sí & 1 & aprueba \\
        \hline
        \texttt{SQL\_injection} & Edge-IIoTset / \texttt{network\_flow} & 7/7 & sí & 1 & aprueba \\
        \hline
        \texttt{Uploading} & Edge-IIoTset / \texttt{network\_packet} & 7/7 & sí & 0 & aprueba \\
        \hline
        \texttt{Vulnerability\_scanner} & Edge-IIoTset / \texttt{network\_packet} & 7/7 & sí & 0 & aprueba \\
        \hline
        \texttt{XSS} & TON-IoT/red / \texttt{network\_flow} & 7/7 & sí & 0 & aprueba \\
        \hline
        \texttt{DoS} & BoT-IoT / \texttt{network\_flow} & 7/7 & sí & 0 & aprueba \\
        \hline
        \texttt{Command\_and\_Control} & IoT-23 / \texttt{network\_flow} & 8/8 & sí & 0 & aprueba \\
        \hline
    \end{tabular}
\end{table}

Las 16 respuestas crudas fijaron
\texttt{requires\_human\_review=true}. El juez neutralizó esta petición en 15
casos porque las cinco primeras recomendaciones cumplían la regla estructural
y no existía otra incidencia. El único caso derivado fue
\texttt{DDoS\_UDP}: cinco bases quedaron contextualizadas, dos acciones
omitidas por el modelo se conservaron como catálogo, se añadieron dos
sugerencias y apareció la referencia \texttt{T1498.003}, que no pertenecía al
catálogo exacto autorizado para ese tipo. El juez generó
\texttt{human\_interrupt} y el auditor emitió \texttt{review}; este último
resultado no es un fallo del auditor, sino la confirmación de que la derivación
se realizó de acuerdo con la política.

Como medida léxica auxiliar se comparó cada uno de los 111 campos
\texttt{context} anclados con la medida base a la que el propio modelo los
asoció. La media de ROUGE-L fue 0,2057 y la similitud coseno TF--IDF fue 0,2813;
ningún contexto fue idéntico a su base tras la normalización. Estos valores
indican que hubo una ampliación textual, pero no miden corrección técnica ni
fidelidad semántica. La base catalogada no es una referencia independiente y
la asociación mediante \texttt{base\_id} procede de la propia respuesta del
modelo.

\subsection*{Inspección cualitativa asistida de las contextualizaciones}

Una inspección cualitativa asistida por IA matizó los resultados estructurales.
No participaron expertos humanos independientes, por lo que esta revisión se
utiliza como cribado de posibles errores y no como una métrica validada de
calidad semántica. Los contextos de \texttt{SQL\_injection} y
\texttt{Vulnerability\_scanner} fueron los mejor respaldados por señales
explícitas: el primero contenía una consulta HTTP con \texttt{SLEEP(5)} y el
segundo un parámetro \texttt{basepath} dirigido a un recurso RFI. Incluso en el
primer caso apareció un error menor, pues la igualdad \texttt{1274=1274} se
describió como falsa.

En siete casos ---\texttt{Backdoor}, \texttt{DDoS\_HTTP},
\texttt{DDoS\_ICMP}, \texttt{DDoS\_TCP}, \texttt{Fingerprinting},
\texttt{Port\_Scanning} y \texttt{Ransomware}--- la redacción fue compatible
con el tipo predicho, pero atribuyó a un único paquete o flujo propiedades no
observables directamente, como persistencia, inundación, múltiples intentos o
propagación. Otros siete casos ---\texttt{DDoS\_UDP}, \texttt{MITM},
\texttt{Password}, \texttt{Uploading}, \texttt{XSS}, \texttt{DoS} y
\texttt{Command\_and\_Control}--- fueron marcados por contener al menos una
posible contradicción o incorrección técnica material. Esta clasificación es
descriptiva y no puede extrapolarse más allá de los 16 ejemplos.

El caso más evidente fue \texttt{DDoS\_UDP}. El evento contenía un SYN TCP y
no identificaba un destino válido, pero Mistral supuso tráfico UDP encapsulado
en TCP, trató el puerto TCP 6476 como si fuese UDP y añadió
\texttt{T1498.003}. También se detectaron errores que la validación estructural
no capturó: en \texttt{DoS}, ocho paquetes en 24,98 segundos
---aproximadamente 0,32 paquetes por segundo--- se interpretaron como
agotamiento sostenido; en \texttt{Command\_and\_Control}, el puerto 50 se
identificó como DNS, cuyo puerto habitual es el 53; en \texttt{Uploading}, una
opción TCP MSS de 1460 bytes y la ausencia normal de ACK en un SYN inicial se
presentaron como posibles indicios de manipulación; y en \texttt{XSS}, sin
\textit{payload} HTTP observable, apareció además el marcador ficticio
\texttt{CVE-2023-XXXX}.

Estos hallazgos no modificaron las acciones operativas: el campo \texttt{text}
continuó conteniendo el literal inmutable del catálogo y la aportación
generativa permaneció separada como contexto no confiable. Sin embargo,
demuestran que la existencia de un \texttt{base\_id} válido no garantiza que
la explicación conserve fielmente la evidencia del evento. Por tanto, las 15
aprobaciones obtenidas deben interpretarse como cumplimiento de las reglas
estructurales del juez, no como 15 contextualizaciones técnicamente correctas.

\subsection{Robustez, trazabilidad y degradación controlada}

La campaña en vivo no observó indisponibilidad del proveedor: las 16 llamadas
fueron procesables y el número de \textit{fallbacks} fue cero. En consecuencia,
esta ejecución no constituye evidencia empírica en vivo de la ruta de reserva.
Dicha ruta se comprobó mediante pruebas automatizadas que simulan una caída del
proveedor para cada uno de los 16 tipos. En esos casos, el mitigador conserva
la respuesta específica del catálogo y registra el intento generativo como
\texttt{fallback}. También se probaron respuestas mal formadas,
\texttt{base\_id} inexistentes o repetidos, confianzas no numéricas y
referencias representadas mediante tipos inesperados.

La suite específica del mitigador superó sus 83 pruebas. Una ejecución conjunta
de las pruebas del mitigador, los agentes finales, el auditor, los contratos del
caso y los servidores MCP superó 292 pruebas, sin fallos funcionales. Los 16
casos de la campaña se recuperaron de la memoria persistente y produjeron 112
entradas de traza, exactamente siete por caso y en el orden esperado. Los 16
\texttt{CaseResult} validaron su contrato; el auditor no registró ningún fallo
duro.

Se conservaron el manifiesto de campaña, los casos seleccionados, las respuestas
JSON de Mistral anteriores al anclaje, los \texttt{CaseResult}, los informes del
auditor, las métricas por caso, el registro HTTP sin cabeceras sensibles, la base
de memoria y los resúmenes. Los 12 ficheros quedaron cubiertos por
\texttt{SHA256SUMS}; todos los hashes se verificaron y los artefactos no contienen
la credencial ni cabeceras de autorización.

\subsection{Alcance y limitaciones de la evaluación}

La campaña está balanceada por tipo únicamente en el sentido de que contiene un
caso de cada una de las 16 clases. No está balanceada por procedencia: diez casos
proceden de Edge-IIoTset, cuatro de TON-IoT/red, uno de BoT-IoT y uno de IoT-23.
Tampoco cubre todos los perfiles: nueve eventos son
\texttt{network\_packet} y siete son \texttt{network\_flow}; no hay casos en
vivo de \texttt{host\_metrics}, \texttt{iot\_telemetry},
\texttt{alert\_text} o \texttt{pcap\_ref}.

Además, la selección priorizó deliberadamente los candidatos con mayor confianza
y condicionó la entrada a una detección y clasificación correctas. Por ello,
esta campaña valida cobertura funcional, contrato, comunicación con el proveedor,
anclaje, persistencia y respuesta del juez para el tramo posterior a la
estandarización. Con una observación por tipo no permite estimar el rendimiento
poblacional, la robustez frente a errores de agentes anteriores ni la
generalización por dataset o perfil.

El manifiesto fija el commit de partida y los hashes del catálogo, el esquema,
el prompt, los modelos, la selección y las entradas. No obstante, la ejecución
se realizó sobre un árbol de trabajo con cambios aún no consolidados, por lo que
el commit por sí solo no reconstruye exactamente el código ejecutado. Antes de
considerar estos resultados como evidencia definitiva de una versión publicada,
la campaña debería repetirse desde una revisión limpia y etiquetada.

En conjunto, la campaña demuestra que el mitigador puede consultar al proveedor
real para los 16 tipos, conservar el texto verificable del catálogo, separar las
ampliaciones generativas, derivar referencias ajenas, persistir el caso y dejar
una traza auditable. No demuestra todavía que las contextualizaciones sean
semánticamente correctas. Los errores detectados muestran que la regla de las
cinco primeras recomendaciones es suficiente como control estructural, pero no
como validación semántica; esta última requiere una rúbrica independiente y una
revisión por expertos antes de sostener una aptitud productiva.
